\documentclass[11pt]{article}

\usepackage[margin=1in]{geometry}
\usepackage{amsmath,amssymb,mathtools}
\usepackage{longtable,array}

\newcommand{\E}{\mathbb E}
\newcommand{\rank}{\operatorname{rank}}

\begin{document}

\begin{center}
\textbf{From the Static BLM Regression to the Project Model}

Complete proof and numerical example
\end{center}

\section*{Scope and notation}

Bonhomme, Lamadon, and Manresa (2019, p.~704, equation~(1)) give the
following static interactive regression:
\[
Y_{it}
=
a_t(k_{it})
+
b_t(k_{it})\alpha_i^B
+
X_{it}'c_t
+
\varepsilon_{it}.
\tag{BLM}
\]
The superscript \(B\) distinguishes the BLM worker type
\(\alpha_i^B\) from the project's additive worker effect \(\alpha_i\).
We fix one period \(t\), suppress time subscripts, and use \(K\) for the
number of worker types and \(L\) for the number of firm classes.  This
relabels BLM's firm-class index so that worker types and firm classes have
different letters.

\section*{Notation table}

\begin{longtable}{
  >{\raggedright\arraybackslash}p{0.22\linewidth}
  >{\raggedright\arraybackslash}p{0.68\linewidth}}
\hline
Symbol & Meaning \\
\hline
\endfirsthead
\hline
Symbol & Meaning \\
\hline
\endhead
\(t\) & Time-period index in the original BLM panel model. \\
\(T\) & Number of periods in the original BLM panel and the endpoint of a full history. \\
\(i\) & Individual worker index, \(i=1,\ldots,I\). \\
\(j\) & Individual firm index, \(j=1,\ldots,J\). \\
\(k\) & Worker-type index, \(k=1,\ldots,K\). \\
\(\ell\) & Firm-class index, \(\ell=1,\ldots,L\). \\
\(I\) & Number of individual workers. \\
\(J\) & Number of individual firms. \\
\(K\) & Number of discrete BLM worker types in the finite-type specialization. \\
\(L\) & Number of BLM firm classes. \\
\(j_{it}\) & Identity of worker \(i\)'s employer in period \(t\). \\
\(k_{it}\) & BLM class of worker \(i\)'s employer in period \(t\). \\
\(m_{it}\) & BLM indicator that worker \(i\) moves firms between periods \(t\) and \(t+1\); unrelated to the conditional mean \(m_{ij}\). \\
\(k_i^T,m_i^T,X_i^T\) & Worker \(i\)'s complete firm-class, mobility, and covariate histories through period \(T\). \\
\(g(i)\) & Worker-type membership map: worker \(i\) has type \(g(i)\). \\
\(h(j)\) & Firm-class membership map: firm \(j\) belongs to class \(h(j)\). \\
\(Y_{it}\) & Worker \(i\)'s observed log earnings in period \(t\). \\
\(Y_{ij}\) & The same outcome after fixing a period and indexing the observation by worker \(i\) and employer \(j\). \\
\(X_{it},X_{ij}\) & Observed wage covariates before and after fixing the period. \\
\(c_t,c\) & BLM covariate coefficients before and after fixing the period. \\
\(\delta\) & Project-model covariate coefficient; the mapping sets \(\delta=c\). \\
\(\varepsilon_{it},\varepsilon_{ij}\) & Mean-zero wage residual before and after fixing the period. \\
\(\alpha_i^B\) & BLM latent type of worker \(i\). \\
\(\alpha_k^B\) & Numerical value of BLM worker type \(k\). \\
\(\alpha_i\) & Project-model additive worker effect; this is not the BLM type. \\
\(\alpha\) & The \(I\times1\) vector of project additive worker effects. \\
\(a_t(k),a_\ell\) & BLM firm-class intercept before and after fixing the period. \\
\(b_t(k),b_\ell\) & BLM firm-class loading on the BLM worker type before and after fixing the period. \\
\(p_k\) & Reference probability of worker type \(k\), with \(\sum_k p_k=1\). \\
\(q_\ell\) & Reference probability of firm class \(\ell\), with \(\sum_\ell q_\ell=1\). \\
\(p\) & The \(K\times1\) vector with entries \(p_k\). \\
\(q\) & The \(L\times1\) vector with entries \(q_\ell\). \\
\(\bar\alpha\) & Reference mean of the BLM worker type, \(\sum_k p_k\alpha_k^B\). \\
\(\bar a\) & Reference mean of the BLM firm-class intercept, \(\sum_\ell q_\ell a_\ell\). \\
\(\bar b\) & Reference mean of the BLM firm-class loading, \(\sum_\ell q_\ell b_\ell\). \\
\(m_{ij}\) & Conditional mean wage for worker \(i\) at firm \(j\). \\
\(\mu_{k\ell}\) & Conditional mean net of covariates for worker type \(k\) at firm class \(\ell\). \\
\(\mu_{k\cdot}\) & Mean of row \(k\), averaging \(\mu_{k\ell}\) over firm classes using \(q_\ell\). \\
\(\mu_{\cdot\ell}\) & Mean of column \(\ell\), averaging \(\mu_{k\ell}\) over worker types using \(p_k\). \\
\(\mu_{\cdot\cdot}\) & Grand mean of the type-class wage table using weights \(p_kq_\ell\). \\
\(\mu\) & Project-model intercept; the mapping sets \(\mu=\mu_{\cdot\cdot}\). \\
\(\psi_j\) & Project-model additive firm effect. \\
\(\psi_\ell\) & Project additive effect shared by every firm in class \(\ell\). \\
\(\psi\) & The \(J\times1\) vector of project additive firm effects. \\
\(M_{k\ell}\) & Double-demeaned interaction between worker type \(k\) and firm class \(\ell\). \\
\(M_{\cdot\ell}\) & Column \(\ell\) of \(M\). \\
\(M\) & The \(K\times L\) matrix with entries \(M_{k\ell}\). \\
\(u_i\) & Project worker interaction factor, \(\alpha_i^B-\bar\alpha\). \\
\(v_j\) & Project firm interaction factor, \(b_{h(j)}-\bar b\). \\
\(u\) & The \(I\times1\) vector with entries \(u_i\). \\
\(v\) & The \(J\times1\) vector with entries \(v_j\). \\
\(\Lambda\) & Project interaction coefficient; in the rank-one BLM mapping, \(\Lambda=[1]\). \\
\(Z_{ik}\) & Worker-type indicator, \(\mathbf 1\{g(i)=k\}\). \\
\(Z\) & The \(I\times K\) worker-type membership matrix. \\
\(W_{j\ell}\) & Firm-class indicator, \(\mathbf 1\{h(j)=\ell\}\). \\
\(W\) & The \(J\times L\) firm-class membership matrix. \\
\(H_{ij}\) & Individual worker-firm interaction, \(M_{g(i),h(j)}\). \\
\(H\) & The \(I\times J\) matrix with entries \(H_{ij}\). \\
\(x\) & Generic \(K\times1\) vector used in the centered subspace \(\{x:p'x=0\}\). \\
\(A\) & Generic matrix used only in the notation \(\rank(A)\). \\
\('\) & Matrix or vector transpose. \\
\(\E[\cdot]\) & Expectation. \\
\(\rank(A)\) & Dimension of the column space of matrix \(A\). \\
\hline
\end{longtable}

\section*{Complete linear-chain proof}

\paragraph{Step 1. Take the conditional mean of the BLM equation.}

After fixing \(t\), write the BLM equation as
\[
Y_{ij}
=
X_{ij}'c
+
a_{h(j)}
+
b_{h(j)}\alpha_i^B
+
\varepsilon_{ij}.
\tag{1}
\]
BLM impose
\[
\E[\varepsilon_{it}\mid
\alpha_i^B,k_i^T,m_i^T,X_i^T]=0.
\tag{2}
\]
The current variables \(X_{ij}\), \(\alpha_i^B\), and \(h(j)=k_{it}\) are
contained in that full conditioning information.  The law of iterated
expectations therefore gives
\begin{align*}
\E[\varepsilon_{ij}\mid X_{ij},\alpha_i^B,h(j)]
&=
\E\!\left[
\E[\varepsilon_{it}\mid\alpha_i^B,k_i^T,m_i^T,X_i^T]
\mid X_{ij},\alpha_i^B,h(j)
\right]\\
&=\E[0\mid X_{ij},\alpha_i^B,h(j)]\\
&=0.
\tag{2a}
\end{align*}
Take the conditional expectation of every term in (1):
\begin{align*}
m_{ij}
&\equiv
\E[Y_{ij}\mid X_{ij},\alpha_i^B,h(j)]\\
&=
X_{ij}'c
+
a_{h(j)}
+
b_{h(j)}\alpha_i^B
+
\E[\varepsilon_{ij}\mid X_{ij},\alpha_i^B,h(j)]\\
&=
X_{ij}'c+a_{h(j)}+b_{h(j)}\alpha_i^B.
\end{align*}
Thus
\[
m_{ij}=X_{ij}'c+a_{h(j)}+b_{h(j)}\alpha_i^B.
\tag{3}
\]

\paragraph{Step 2. Write the mean at the type-class level.}

Suppose worker \(i\) has type \(g(i)\), so
\[
\alpha_i^B=\alpha_{g(i)}^B.
\tag{4}
\]
For worker type \(k\) and firm class \(\ell\), define
\[
\mu_{k\ell}=a_\ell+b_\ell\alpha_k^B.
\tag{5}
\]
Set \(k=g(i)\) and \(\ell=h(j)\) in (5):
\[
\mu_{g(i),h(j)}
=
a_{h(j)}+b_{h(j)}\alpha_{g(i)}^B
=
a_{h(j)}+b_{h(j)}\alpha_i^B.
\tag{6}
\]
Substituting (6) into (3) gives
\[
m_{ij}=X_{ij}'c+\mu_{g(i),h(j)}.
\tag{7}
\]

\paragraph{Step 3. Define the reference means.}

Let
\[
\sum_{k=1}^Kp_k=1,
\qquad
\sum_{\ell=1}^Lq_\ell=1.
\tag{8}
\]
Define
\[
\bar\alpha=\sum_{k=1}^Kp_k\alpha_k^B,
\qquad
\bar a=\sum_{\ell=1}^Lq_\ell a_\ell,
\qquad
\bar b=\sum_{\ell=1}^Lq_\ell b_\ell.
\tag{9}
\]

\paragraph{Step 4. Calculate each row mean.}

For worker type \(k\),
\begin{align*}
\mu_{k\cdot}
&\equiv
\sum_{\ell=1}^Lq_\ell\mu_{k\ell}\\
&=
\sum_{\ell=1}^Lq_\ell(a_\ell+b_\ell\alpha_k^B)\\
&=
\sum_{\ell=1}^Lq_\ell a_\ell
+
\alpha_k^B\sum_{\ell=1}^Lq_\ell b_\ell\\
&=
\bar a+\bar b\alpha_k^B.
\end{align*}
Therefore,
\[
\mu_{k\cdot}=\bar a+\bar b\alpha_k^B.
\tag{10}
\]

\paragraph{Step 5. Calculate each column mean.}

For firm class \(\ell\),
\begin{align*}
\mu_{\cdot\ell}
&\equiv
\sum_{k=1}^Kp_k\mu_{k\ell}\\
&=
\sum_{k=1}^Kp_k(a_\ell+b_\ell\alpha_k^B)\\
&=
a_\ell\sum_{k=1}^Kp_k
+
b_\ell\sum_{k=1}^Kp_k\alpha_k^B\\
&=
a_\ell+b_\ell\bar\alpha.
\end{align*}
Therefore,
\[
\mu_{\cdot\ell}=a_\ell+b_\ell\bar\alpha.
\tag{11}
\]

\paragraph{Step 6. Calculate the grand mean.}

Define
\[
\mu_{\cdot\cdot}
\equiv
\sum_{k=1}^K\sum_{\ell=1}^Lp_kq_\ell\mu_{k\ell}.
\tag{12}
\]
Substitute (5):
\begin{align*}
\mu_{\cdot\cdot}
&=
\sum_{k=1}^K\sum_{\ell=1}^L
p_kq_\ell(a_\ell+b_\ell\alpha_k^B)\\
&=
\left(\sum_{k=1}^Kp_k\right)
\left(\sum_{\ell=1}^Lq_\ell a_\ell\right)
+
\left(\sum_{k=1}^Kp_k\alpha_k^B\right)
\left(\sum_{\ell=1}^Lq_\ell b_\ell\right)\\
&=
1\cdot\bar a+\bar\alpha\bar b.
\end{align*}
Hence
\[
\mu_{\cdot\cdot}=\bar a+\bar\alpha\bar b.
\tag{13}
\]
The project intercept is
\[
\mu=\mu_{\cdot\cdot}=\bar a+\bar\alpha\bar b.
\tag{14}
\]

\paragraph{Step 7. Obtain the project additive worker effect.}

The additive worker effect is the worker's row mean minus the grand mean:
\[
\alpha_i
\equiv
\mu_{g(i),\cdot}-\mu_{\cdot\cdot}.
\tag{15}
\]
Using (10), (13), and \(\alpha_{g(i)}^B=\alpha_i^B\),
\begin{align*}
\alpha_i
&=
\bar a+\bar b\alpha_i^B
-
(\bar a+\bar\alpha\bar b)\\
&=
\bar b\alpha_i^B-\bar b\bar\alpha\\
&=
\bar b(\alpha_i^B-\bar\alpha).
\end{align*}
Therefore,
\[
\alpha_i=\bar b(\alpha_i^B-\bar\alpha).
\tag{16}
\]

\paragraph{Step 8. Obtain the project additive firm effect.}

The additive firm effect is the firm's column mean minus the grand mean:
\[
\psi_j
\equiv
\mu_{\cdot,h(j)}-\mu_{\cdot\cdot}.
\tag{17}
\]
Using (11) and (13),
\begin{align*}
\psi_j
&=
a_{h(j)}+b_{h(j)}\bar\alpha
-
(\bar a+\bar\alpha\bar b)\\
&=
(a_{h(j)}-\bar a)
+
\bar\alpha(b_{h(j)}-\bar b).
\end{align*}
Therefore,
\[
\psi_j
=
(a_{h(j)}-\bar a)
+
\bar\alpha(b_{h(j)}-\bar b).
\tag{18}
\]

\paragraph{Step 9. Define and calculate the interaction matrix \(M\).}

Double demeaning defines
\[
M_{k\ell}
=
\mu_{k\ell}
-
\mu_{k\cdot}
-
\mu_{\cdot\ell}
+
\mu_{\cdot\cdot}.
\tag{19}
\]
Substitute (5), (10), (11), and (13):
\begin{align*}
M_{k\ell}
&=
(a_\ell+b_\ell\alpha_k^B)
-
(\bar a+\bar b\alpha_k^B)\\
&\quad
-
(a_\ell+b_\ell\bar\alpha)
+
(\bar a+\bar\alpha\bar b)\\
&=
a_\ell+b_\ell\alpha_k^B
-
\bar a-\bar b\alpha_k^B
-
a_\ell-b_\ell\bar\alpha
+
\bar a+\bar\alpha\bar b\\
&=
b_\ell\alpha_k^B
-
\bar b\alpha_k^B
-
b_\ell\bar\alpha
+
\bar\alpha\bar b\\
&=
\alpha_k^B(b_\ell-\bar b)
-
\bar\alpha(b_\ell-\bar b)\\
&=
(\alpha_k^B-\bar\alpha)(b_\ell-\bar b).
\end{align*}
Thus
\[
M_{k\ell}
=
(\alpha_k^B-\bar\alpha)(b_\ell-\bar b).
\tag{20}
\]

\paragraph{Step 10. Prove the rank of \(M\).}

Writing all entries of \(M\) gives
\[
M
=
\begin{pmatrix}
\alpha_1^B-\bar\alpha\\
\vdots\\
\alpha_K^B-\bar\alpha
\end{pmatrix}
\begin{pmatrix}
b_1-\bar b&\cdots&b_L-\bar b
\end{pmatrix}.
\tag{21}
\]
For every \(\ell\), column \(\ell\) is
\[
M_{\cdot\ell}
=
(b_\ell-\bar b)
\begin{pmatrix}
\alpha_1^B-\bar\alpha\\
\vdots\\
\alpha_K^B-\bar\alpha
\end{pmatrix}.
\tag{22}
\]
Every column is therefore a scalar multiple of one vector.  The column space
has dimension at most one:
\[
\rank(M)\leq1.
\tag{23}
\]
This conclusion is obtained before \(H\) is introduced.

\paragraph{Step 11. Prove the general double-centering bound.}

For every worker type \(k\),
\begin{align*}
\sum_{\ell=1}^Lq_\ell M_{k\ell}
&=
(\alpha_k^B-\bar\alpha)
\sum_{\ell=1}^Lq_\ell(b_\ell-\bar b)\\
&=
(\alpha_k^B-\bar\alpha)
\left(
\sum_{\ell=1}^Lq_\ell b_\ell
-
\bar b\sum_{\ell=1}^Lq_\ell
\right)\\
&=
(\alpha_k^B-\bar\alpha)(\bar b-\bar b)\\
&=0.
\end{align*}
Hence
\[
Mq=0.
\tag{24}
\]
Because \(q\neq0\), the right null space of \(M\) has dimension at least one.
Rank--nullity gives
\[
\rank(M)\leq L-1.
\tag{25}
\]
Similarly, for every firm class \(\ell\),
\begin{align*}
\sum_{k=1}^Kp_kM_{k\ell}
&=
(b_\ell-\bar b)
\sum_{k=1}^Kp_k(\alpha_k^B-\bar\alpha)\\
&=
(b_\ell-\bar b)
\left(
\sum_{k=1}^Kp_k\alpha_k^B
-
\bar\alpha\sum_{k=1}^Kp_k
\right)\\
&=
(b_\ell-\bar b)(\bar\alpha-\bar\alpha)\\
&=0.
\end{align*}
Hence
\[
p'M=0.
\tag{26}
\]
The columns of \(M\) lie in the \((K-1)\)-dimensional subspace
\(\{x\in\mathbb R^K:p'x=0\}\), so
\[
\rank(M)\leq K-1.
\tag{27}
\]
Combining (23), (25), and (27),
\[
\rank(M)\leq\min\{1,K-1,L-1\}.
\tag{28}
\]

\paragraph{Step 12. Recover the type-class decomposition.}

Starting from (19),
\[
M_{k\ell}
=
\mu_{k\ell}-\mu_{k\cdot}-\mu_{\cdot\ell}+\mu_{\cdot\cdot}.
\]
Add \(\mu_{k\cdot}+\mu_{\cdot\ell}-\mu_{\cdot\cdot}\) to both sides:
\[
\mu_{k\ell}
=
\mu_{k\cdot}+\mu_{\cdot\ell}-\mu_{\cdot\cdot}+M_{k\ell}.
\tag{29}
\]
Rewrite the right side by adding and subtracting one more grand mean:
\[
\mu_{k\ell}
=
\mu_{\cdot\cdot}
+
(\mu_{k\cdot}-\mu_{\cdot\cdot})
+
(\mu_{\cdot\ell}-\mu_{\cdot\cdot})
+
M_{k\ell}.
\tag{30}
\]

\paragraph{Step 13. Define the membership matrices.}

Define
\[
Z_{ik}
=
\begin{cases}
1,&g(i)=k,\\
0,&g(i)\neq k,
\end{cases}
\qquad
W_{j\ell}
=
\begin{cases}
1,&h(j)=\ell,\\
0,&h(j)\neq\ell.
\end{cases}
\tag{31}
\]
Each row of \(Z\) contains exactly one \(1\), and each row of \(W\) contains
exactly one \(1\).

\paragraph{Step 14. Define the individual interaction and prove \(H=ZMW'\).}

The type-class interaction faced by worker \(i\) at firm \(j\) is
\[
H_{ij}=M_{g(i),h(j)}.
\tag{32}
\]
Now calculate the \((i,j)\) entry of \(ZMW'\).  First,
\[
(ZMW')_{ij}
=
\sum_{\ell=1}^L(ZM)_{i\ell}W'_{\ell j}.
\tag{33}
\]
Second,
\[
(ZM)_{i\ell}
=
\sum_{k=1}^KZ_{ik}M_{k\ell}.
\tag{34}
\]
Substitute (34) into (33) and use \(W'_{\ell j}=W_{j\ell}\):
\[
(ZMW')_{ij}
=
\sum_{\ell=1}^L\sum_{k=1}^K
Z_{ik}M_{k\ell}W_{j\ell}.
\tag{35}
\]
For fixed \(i\), \(Z_{ik}=0\) except when \(k=g(i)\).  Therefore,
\[
\sum_{k=1}^KZ_{ik}M_{k\ell}
=
M_{g(i),\ell}.
\tag{36}
\]
Use (36) in (35):
\[
(ZMW')_{ij}
=
\sum_{\ell=1}^LM_{g(i),\ell}W_{j\ell}.
\tag{37}
\]
For fixed \(j\), \(W_{j\ell}=0\) except when \(\ell=h(j)\).  Therefore,
\[
\sum_{\ell=1}^LM_{g(i),\ell}W_{j\ell}
=
M_{g(i),h(j)}.
\tag{38}
\]
Equations (32), (37), and (38) imply
\[
(ZMW')_{ij}=H_{ij}
\]
for every \(i,j\).  Hence
\[
H=ZMW'.
\tag{39}
\]
No rank assumption is used in this selector argument.

\paragraph{Step 15. Connect \(H\) to the established project factors.}

The project interaction factors are
\[
u_i=\alpha_i^B-\bar\alpha,
\qquad
v_j=b_{h(j)}-\bar b.
\tag{40}
\]
Using (20) and (32),
\begin{align*}
H_{ij}
&=
M_{g(i),h(j)}\\
&=
(\alpha_{g(i)}^B-\bar\alpha)
(b_{h(j)}-\bar b)\\
&=
(\alpha_i^B-\bar\alpha)
(b_{h(j)}-\bar b)\\
&=
u_iv_j.
\end{align*}
Therefore,
\[
H_{ij}=u_iv_j.
\tag{41}
\]
The \((i,j)\) entry of \(uv'\) is \(u_iv_j\), so
\[
H=uv'.
\tag{42}
\]
Combining (39) and (42),
\[
H=ZMW'=uv'.
\tag{43}
\]
Because every column of \(uv'\) is a scalar multiple of \(u\),
\[
\rank(H)\leq1.
\tag{44}
\]
This is a conclusion, not an assumption.

\paragraph{Step 16. Arrive at the project model.}

Set \(k=g(i)\) and \(\ell=h(j)\) in (30):
\begin{align*}
\mu_{g(i),h(j)}
&=
\mu_{\cdot\cdot}
+
(\mu_{g(i),\cdot}-\mu_{\cdot\cdot})\\
&\quad+
(\mu_{\cdot,h(j)}-\mu_{\cdot\cdot})
+
M_{g(i),h(j)}\\
&=
\mu+\alpha_i+\psi_j+H_{ij}\\
&=
\mu+\alpha_i+\psi_j+u_iv_j.
\end{align*}
Substitute this expression into (7):
\[
m_{ij}
=
X_{ij}'c+\mu+\alpha_i+\psi_j+u_iv_j.
\tag{45}
\]
Set
\[
\delta=c,
\qquad
\Lambda=[1].
\tag{46}
\]
Because \(u_i'\Lambda v_j=u_iv_j\) for scalar factors,
\[
m_{ij}
=
X_{ij}'\delta
+
\mu
+
\alpha_i
+
\psi_j
+
u_i'\Lambda v_j.
\tag{Project}
\]
This is the project model.

\paragraph{Step 17. Verify the map by direct substitution.}

The map is
\[
\begin{aligned}
\delta&=c,\\
\mu&=\bar a+\bar\alpha\bar b,\\
\alpha_i&=\bar b(\alpha_i^B-\bar\alpha),\\
\psi_j&=(a_{h(j)}-\bar a)
       +\bar\alpha(b_{h(j)}-\bar b),\\
u_i&=\alpha_i^B-\bar\alpha,\\
v_j&=b_{h(j)}-\bar b,\\
\Lambda&=[1].
\end{aligned}
\tag{47}
\]
Substitute (47) into the project model:
\begin{align*}
m_{ij}
&=
X_{ij}'c
+
(\bar a+\bar\alpha\bar b)
+
\bar b(\alpha_i^B-\bar\alpha)\\
&\quad+
(a_{h(j)}-\bar a)
+
\bar\alpha(b_{h(j)}-\bar b)\\
&\quad+
(\alpha_i^B-\bar\alpha)(b_{h(j)}-\bar b).
\end{align*}
Expand:
\begin{align*}
m_{ij}
&=
X_{ij}'c
+
\bar a+\bar\alpha\bar b
+
\bar b\alpha_i^B-\bar b\bar\alpha\\
&\quad+
a_{h(j)}-\bar a
+
\bar\alpha b_{h(j)}-\bar\alpha\bar b\\
&\quad+
\alpha_i^Bb_{h(j)}
-
\alpha_i^B\bar b
-
\bar\alpha b_{h(j)}
+
\bar\alpha\bar b.
\end{align*}
Cancel equal positive and negative terms:
\[
\bar a-\bar a=0,
\qquad
\bar b\alpha_i^B-\alpha_i^B\bar b=0,
\]
\[
\bar\alpha b_{h(j)}-\bar\alpha b_{h(j)}=0,
\]
and the remaining \(\bar\alpha\bar b\) terms also sum to zero.  Therefore,
\[
m_{ij}
=
X_{ij}'c+a_{h(j)}+b_{h(j)}\alpha_i^B,

\]
which is exactly the BLM conditional mean in (3).

\section*{Numerical example: Stefan, Michael, GS, Meta, and BoA}

\paragraph{Step 1. Workers, firms, types, and classes.}

There are two workers:
\[
I=2,
\qquad
i=1\text{ is Stefan},
\qquad
i=2\text{ is Michael}.
\]
There are three firms:
\[
J=3,
\qquad
j=1\text{ is GS},
\quad
j=2\text{ is Meta},
\quad
j=3\text{ is BoA}.
\]
There are two worker types and two firm classes:
\[
K=2,
\qquad
L=2.
\]
The worker-type assignments are
\[
g(\text{Stefan})=1,
\qquad
g(\text{Michael})=2.
\tag{48}
\]
The firm-class assignments are
\[
h(\text{GS})=1,
\qquad
h(\text{Meta})=2,
\qquad
h(\text{BoA})=1.
\tag{49}
\]

\paragraph{Step 2. Choose numerical BLM parameters.}

Set the BLM worker-type values to
\[
\alpha_{\text{Stefan}}^B=\alpha_1^B=1,
\qquad
\alpha_{\text{Michael}}^B=\alpha_2^B=-1.
\tag{50}
\]
Set the BLM class intercepts and slopes to
\[
a_1=4,
\qquad
a_2=5,
\qquad
b_1=2,
\qquad
b_2=-1.
\tag{51}
\]
Set \(X_{ij}'c=0\) for every worker-firm pair.  The mean wage is therefore
\[
m_{ij}=a_{h(j)}+b_{h(j)}\alpha_i^B.
\tag{52}
\]

\paragraph{Step 3. Calculate every BLM type-class mean.}

For type 1 at class 1,
\[
\mu_{11}=a_1+b_1\alpha_1^B=4+2(1)=6.
\]
For type 1 at class 2,
\[
\mu_{12}=a_2+b_2\alpha_1^B=5+(-1)(1)=4.
\]
For type 2 at class 1,
\[
\mu_{21}=a_1+b_1\alpha_2^B=4+2(-1)=2.
\]
For type 2 at class 2,
\[
\mu_{22}=a_2+b_2\alpha_2^B=5+(-1)(-1)=6.
\]
Thus
\[
(\mu_{k\ell})
=
\begin{pmatrix}
6&4\\
2&6
\end{pmatrix}.
\tag{53}
\]

\paragraph{Step 4. Choose and interpret the reference weights.}

Give Stefan and Michael equal worker weight:
\[
p_1=p_2=\frac12.
\tag{54}
\]
Two of the three firms are in class 1, while one is in class 2:
\[
q_1=\frac23,
\qquad
q_2=\frac13.
\tag{55}
\]
Therefore,
\[
\bar\alpha
=
\frac12(1)+\frac12(-1)
=0,
\tag{56}
\]
\[
\bar a
=
\frac23(4)+\frac13(5)
=
\frac{13}{3},
\tag{57}
\]
and
\[
\bar b
=
\frac23(2)+\frac13(-1)
=1.
\tag{58}
\]

\paragraph{Step 5. Calculate row, column, and grand means.}

Stefan's row mean is
\[
\mu_{1\cdot}
=
\frac23(6)+\frac13(4)
=
\frac{16}{3}.
\tag{59}
\]
Michael's row mean is
\[
\mu_{2\cdot}
=
\frac23(2)+\frac13(6)
=
\frac{10}{3}.
\tag{60}
\]
The class-1 column mean is
\[
\mu_{\cdot1}
=
\frac12(6)+\frac12(2)
=4.
\tag{61}
\]
The class-2 column mean is
\[
\mu_{\cdot2}
=
\frac12(4)+\frac12(6)
=5.
\tag{62}
\]
The grand mean is
\[
\mu_{\cdot\cdot}
=
\frac12\left(\frac{16}{3}\right)
+
\frac12\left(\frac{10}{3}\right)
=
\frac{13}{3}.
\tag{63}
\]
Hence
\[
\mu=\frac{13}{3}.
\tag{64}
\]

\paragraph{Step 6. Calculate the project additive effects.}

For Stefan,
\[
\alpha_{\text{Stefan}}
=
\mu_{1\cdot}-\mu_{\cdot\cdot}
=
\frac{16}{3}-\frac{13}{3}
=1.
\tag{65}
\]
For Michael,
\[
\alpha_{\text{Michael}}
=
\mu_{2\cdot}-\mu_{\cdot\cdot}
=
\frac{10}{3}-\frac{13}{3}
=-1.
\tag{66}
\]
For a class-1 firm,
\[
\psi_1
=
\mu_{\cdot1}-\mu_{\cdot\cdot}
=
4-\frac{13}{3}
=-\frac13.
\tag{67}
\]
For a class-2 firm,
\[
\psi_2
=
\mu_{\cdot2}-\mu_{\cdot\cdot}
=
5-\frac{13}{3}
=\frac23.
\tag{68}
\]
Therefore,
\[
\alpha
=
\begin{pmatrix}
1\\
-1
\end{pmatrix},
\qquad
\psi
=
\begin{pmatrix}
-1/3\\
2/3\\
-1/3
\end{pmatrix},
\tag{69}
\]
where the entries of \(\psi\) are ordered as GS, Meta, and BoA.

\paragraph{Step 7. Calculate every entry of \(M\).}

\begin{align*}
M_{11}
&=
6-\frac{16}{3}-4+\frac{13}{3}
=1,\\
M_{12}
&=
4-\frac{16}{3}-5+\frac{13}{3}
=-2,\\
M_{21}
&=
2-\frac{10}{3}-4+\frac{13}{3}
=-1,\\
M_{22}
&=
6-\frac{10}{3}-5+\frac{13}{3}
=2.
\end{align*}
Thus
\[
M
=
\begin{pmatrix}
1&-2\\
-1&2
\end{pmatrix}.
\tag{70}
\]
It factors as
\[
M
=
\begin{pmatrix}
1\\
-1
\end{pmatrix}
\begin{pmatrix}
1&-2
\end{pmatrix},
\tag{71}
\]
so
\[
\rank(M)=1.
\tag{72}
\]

\paragraph{Step 8. Write the membership matrices.}

Stefan is type 1 and Michael is type 2, so
\[
Z
=
\begin{pmatrix}
1&0\\
0&1
\end{pmatrix}.
\tag{73}
\]
GS, Meta, and BoA have classes 1, 2, and 1, so
\[
W
=
\begin{pmatrix}
1&0\\
0&1\\
1&0
\end{pmatrix},
\qquad
W'
=
\begin{pmatrix}
1&0&1\\
0&1&0
\end{pmatrix}.
\tag{74}
\]

\paragraph{Step 9. Calculate the established project interaction factors.}

Using \(u_i=\alpha_i^B-\bar\alpha\),
\[
u_{\text{Stefan}}=1-0=1,
\qquad
u_{\text{Michael}}=-1-0=-1.
\]
Therefore,
\[
u=
\begin{pmatrix}
1\\
-1
\end{pmatrix}.
\tag{75}
\]
Using \(v_j=b_{h(j)}-\bar b\),
\[
v_{\text{GS}}=2-1=1,
\]
\[
v_{\text{Meta}}=-1-1=-2,
\]
\[
v_{\text{BoA}}=2-1=1.
\]
Therefore,
\[
v=
\begin{pmatrix}
1\\
-2\\
1
\end{pmatrix}.
\tag{76}
\]

\paragraph{Step 10. Calculate \(H\) in both ways.}

First use \(H=ZMW'\):
\begin{align*}
H
&=
\begin{pmatrix}
1&0\\
0&1
\end{pmatrix}
\begin{pmatrix}
1&-2\\
-1&2
\end{pmatrix}
\begin{pmatrix}
1&0&1\\
0&1&0
\end{pmatrix}\\
&=

\begin{pmatrix}
1&-2&1\\
-1&2&-1
\end{pmatrix}.

\end{align*}
Second use \(H=uv'\):
\begin{align*}
H
&=
\begin{pmatrix}
1\\
-1
\end{pmatrix}
\begin{pmatrix}
1&-2&1
\end{pmatrix}\\
&=

\begin{pmatrix}
1&-2&1\\
-1&2&-1
\end{pmatrix}.

\end{align*}
The rows are Stefan and Michael.  The columns are GS, Meta, and BoA.

\paragraph{Step 11. Reconstruct all six worker-firm means.}

The project model is
\[
m_{ij}=\mu+\alpha_i+\psi_j+u_iv_j.
\tag{77}
\]
For Stefan at GS,
\[
m_{\text{Stefan,GS}}
=
\frac{13}{3}+1-\frac13+(1)(1)
=6.
\]
For Stefan at Meta,
\[
m_{\text{Stefan,Meta}}
=
\frac{13}{3}+1+\frac23+(1)(-2)
=4.
\]
For Stefan at BoA,
\[
m_{\text{Stefan,BoA}}
=
\frac{13}{3}+1-\frac13+(1)(1)
=6.
\]
For Michael at GS,
\[
m_{\text{Michael,GS}}
=
\frac{13}{3}-1-\frac13+(-1)(1)
=2.
\]
For Michael at Meta,
\[
m_{\text{Michael,Meta}}
=
\frac{13}{3}-1+\frac23+(-1)(-2)
=6.
\]
For Michael at BoA,
\[
m_{\text{Michael,BoA}}
=
\frac{13}{3}-1-\frac13+(-1)(1)
=2.
\]
Therefore,
\[
(m_{ij})
=
\begin{pmatrix}
6&4&6\\
2&6&2
\end{pmatrix}.
\tag{78}
\]
The same matrix follows directly from the BLM equation
\(m_{ij}=a_{h(j)}+b_{h(j)}\alpha_i^B\).

\paragraph{Step 12. Numerical interpretation of every model component.}

\begin{center}
\begin{tabular}{
  >{\raggedright\arraybackslash}p{0.22\linewidth}
  >{\raggedright\arraybackslash}p{0.22\linewidth}
  >{\raggedright\arraybackslash}p{0.46\linewidth}}
\hline
Object & Numerical value & Interpretation \\
\hline
\(\alpha_{\text{Stefan}}^B,\alpha_{\text{Michael}}^B\)
& \(1,-1\)
& BLM worker types. \\
\(a_1,a_2\)
& \(4,5\)
& Baseline wages for firm classes 1 and 2 when the BLM worker type is zero. \\
\(b_1,b_2\)
& \(2,-1\)
& Class-specific loadings on the BLM worker type. \\
\(p_1,p_2\)
& \(1/2,1/2\)
& Worker-type reference weights. \\
\(q_1,q_2\)
& \(2/3,1/3\)
& Firm-class reference weights. \\
\(\bar\alpha,\bar a,\bar b\)
& \(0,13/3,1\)
& Reference means used for double demeaning. \\
\(\mu\)
& \(13/3\)
& Overall project intercept. \\
\(\alpha_{\text{Stefan}},\alpha_{\text{Michael}}\)
& \(1,-1\)
& Project additive worker effects. \\
\(\psi_{\text{GS}},\psi_{\text{Meta}},\psi_{\text{BoA}}\)
& \(-1/3,2/3,-1/3\)
& Project additive firm effects. \\
\(u_{\text{Stefan}},u_{\text{Michael}}\)
& \(1,-1\)
& Project worker interaction factors. \\
\(v_{\text{GS}},v_{\text{Meta}},v_{\text{BoA}}\)
& \(1,-2,1\)
& Project firm interaction factors. \\
\(M\)
& \(\begin{psmallmatrix}1&-2\\-1&2\end{psmallmatrix}\)
& Interaction at the worker-type/firm-class level. \\
\(H\)
& \(\begin{psmallmatrix}1&-2&1\\-1&2&-1\end{psmallmatrix}\)
& Interaction at the individual worker/firm level. \\
\(\Lambda\)
& \(1\)
& Rank-one project interaction coefficient. \\
\hline
\end{tabular}
\end{center}

\section*{Dependency check}

The proof has the following one-directional dependency structure:
\[
\begin{aligned}
\text{BLM equation and mean-zero residual}
&\Longrightarrow
\text{conditional mean }m_{ij}\\
&\Longrightarrow
\text{type-class means }\mu_{k\ell}\\
&\Longrightarrow
\text{row, column, and grand means}\\
&\Longrightarrow
\text{double-demeaned }M_{k\ell}\\
&\Longrightarrow
M_{k\ell}
=(\alpha_k^B-\bar\alpha)(b_\ell-\bar b)\\
&\Longrightarrow
\rank(M)\leq1\\
&\Longrightarrow
H_{ij}=M_{g(i),h(j)}\\
&\Longrightarrow
H=ZMW'\\
&\Longrightarrow
H_{ij}=u_iv_j\\
&\Longrightarrow
H=uv'\\
&\Longrightarrow
m_{ij}
=X_{ij}'\delta+\mu+\alpha_i+\psi_j+u_i'\Lambda v_j.
\end{aligned}
\]
In particular, the rank of \(H\) is never used to prove the rank of \(M\).
The rank of \(M\) is established first from the BLM coefficients.

\section*{Scope of the result}

The rank-one conclusion applies to BLM's displayed static interactive
regression
\[
a_\ell+b_\ell\alpha_k^B.
\]
BLM's general static framework permits richer conditional earnings
distributions.  For an arbitrary finite \(K\times L\) conditional-mean table,
double demeaning still implies
\[
\rank(M)\leq\min\{K-1,L-1\},
\]
but it need not imply rank one.  The derivation above therefore proves exact
rank-one nesting for BLM's displayed interactive regression and exact
finite-rank nesting for a general finite type-class mean table.  It nests the
static conditional-mean layer, not BLM's complete earnings distribution,
assignment process, or dynamic mobility model.

\begin{thebibliography}{9}
\bibitem{BLM2019}
Bonhomme, S., Lamadon, T., and Manresa, E. (2019).
``A Distributional Framework for Matched Employer Employee Data.''
\emph{Econometrica}, 87(3), 699--739.
\end{thebibliography}

\end{document}
